\section*{Document index}
\addcontentsline{toc}{section}{Document index}
\label{sec:docindex}

Every source file in \texttt{paper/}, what it holds, and how far along it is.
\statusdone{} means real prose exists; \statusscaffold{} means planning notes
only; \statusblocked{} means the section cannot be written until data lands.

\bigskip
\noindent
\begin{tabular}{@{}p{0.30\textwidth} p{0.44\textwidth} l@{}}
\toprule
\textbf{File} & \textbf{Contents} & \textbf{Status} \\
\midrule
\multicolumn{3}{@{}l}{\textit{Build}} \\
\texttt{main.tex}      & Compile order; \textbackslash input's everything     & \statusdone \\
\texttt{preamble.tex}  & Packages, palette, \texttt{\textbackslash plan} /
                         \texttt{\textbackslash todo} macros, draft-mode switch & \statusdone \\
\texttt{Makefile}      & \texttt{make pdf}; SVG$\rightarrow$PDF via Inkscape  & \statusdone \\
\texttt{refs.bib}      & Bibliography; all entries machine-retrieved          & \statusdone \\
\addlinespace
\multicolumn{3}{@{}l}{\textit{Front matter}} \\
\texttt{abstract.tex}  & $\sim$125-word abstract, written last                & \statusscaffold \\
\texttt{intro.tex}     & Four-paragraph introduction arc                      & \statusscaffold \\
\texttt{related\_work.tex} & Status quo: backbone generation, filtering,
                         RL for inverse folding, library design, the gap      & \statusdone \\
\addlinespace
\multicolumn{3}{@{}l}{\textit{Results} --- the three-move spine} \\
\texttt{results/frontier.tex}  & Move 1: RL moves the diversity/fidelity
                         frontier. Public benchmark, referee-reproducible  & \statusscaffold \\
\texttt{results/fragments.tex} & Move 2: split + recombine $\rightarrow$
                         combinatorial library. \textbf{No data yet}        & \statusblocked \\
\texttt{results/petase.tex}    & Move 3: three closed-loop rounds to
                         $250\times$; the demonstration                     & \statusblocked \\
\addlinespace
\multicolumn{3}{@{}l}{\textit{Methods}} \\
\texttt{methods/overview.tex} & GRPO, fragment libraries, surrogate          & \statusscaffold \\
\texttt{methods/petase.tex}   & Three libraries, surrogate, assays           & \statusscaffold \\
\addlinespace
\multicolumn{3}{@{}l}{\textit{Back matter}} \\
\texttt{conclusion.tex} & Four-paragraph discussion and outlook              & \statusscaffold \\
\addlinespace
\multicolumn{3}{@{}l}{\textit{Not in this build}} \\
\texttt{notebook.tex} & Separate document: the E1--E21 record               & \statusdone \\
\texttt{experiments.tex} & Arms, selection rules, slate E1--E21, algorithm
                        audit; \textbf{\textbackslash input by notebook.tex,
                        not by main.tex}                                     & \statusdone \\
\texttt{supplementary/*} & Heme and protease results + methods, cut from the
                        spine and parked                                    & \statusblocked \\
\bottomrule
\end{tabular}

\bigskip
\noindent
\textbf{The paper and the record are two documents.} \texttt{main.tex} builds
the paper (\texttt{make pdf}); \texttt{notebook.tex} builds the experimental
record (\texttt{make notebook}); \texttt{make all} builds both. The split is
deliberate: the notebook keeps pre-registrations next to the results that
falsified them, which is what makes a retraction checkable and what a paper
cannot carry and stay readable. A claim moves from the notebook to the paper by
being rewritten as prose, never by \textbackslash input.

\bigskip
\noindent
Two companion files live in \texttt{paper/} but are deliberately \emph{not}
part of either build:
\begin{itemize}\setlength{\itemsep}{0pt}
  \item \texttt{outline\_computational.md} --- working brainstorm for a
        computational-first version of the paper: landscape analysis, the
        three candidate framings, seven proposed experiments (E1--E7), figure
        plan, and risk register. Section~\ref{sec:related} is the LaTeX
        rendering of its landscape half.
  \item \texttt{figures/palette.py} --- the plotting counterpart to the
        \texttt{cleo*} colors defined in \texttt{preamble.tex}.
\end{itemize}

\bigskip
\noindent
\textbf{Draft mode.} This document currently compiles with
\texttt{\textbackslash draftmodetrue} in \texttt{preamble.tex}, which renders
all planning notes in gray and all \texttt{\textbackslash todo} markers in red.
Setting \texttt{\textbackslash draftmodefalse} hides every one of them, leaving
only submission prose --- useful for checking how much real text actually
exists.
